\documentclass[11pt]{article}
\usepackage[margin=1in]{geometry}
\usepackage{booktabs}
\usepackage{array}
\usepackage{enumitem}
\usepackage[hidelinks]{hyperref}
\setlength{\parskip}{0.5em}
\setlength{\parindent}{0pt}
\renewcommand{\arraystretch}{1.15}

\title{Progress Report -- Week [6]}
\author{%
  \begin{tabular}[t]{c@{\hspace{2em}}c@{\hspace{2em}}c}
    Aniyah Mack & Adonay Zeru & Hasset Admasu\\
    \texttt{amack18@emich.edu} & \texttt{azeru@emich.edu} & \texttt{hadmasu@emich.edu}
  \end{tabular}\\[0.8em]
  Nathnael Abraham\\
  \texttt{nabrah10@emich.edu}}
\date{October 26, 2026}

\begin{document}
\maketitle

\section*{Project Overview}
Our project investigates how vision-language models can identify important events in soccer match
videos and generate natural-language commentary describing those events. Our main goal is to
determine how accurately soccer events can be identified from video, and how those events can be
used to generate accurate commentary for different audiences, including casual viewers, expert fans,
and people with low vision. The project builds on the SoccerReplay-1988 dataset~[5] and existing
baseline methods for soccer event classification.

\section*{Where Our Last Report Left Off}
Last week we cleaned up the pipeline on the GPU server: balanced accuracy in every report, splits by whole
match, and documented label exclusions. We promised baseline results on real data for this report. Your
feedback answered our two questions (keep the rare classes separate until their counts and per-class
scores are known; hold out the most recent unused season) and asked for baseline results with macro-F1
and balanced accuracy before more experiments. This report delivers those results and acts on both answers.
Everything here is on our \texttt{cleanup/dino-contrastive} branch, and a summary is in our weekly check-in
repository~[7].

\section*{Current Progress}

\subsection*{1. Data: 255 Premier League matches}
We extended the data from 65 matches (2021--22) to 255 matches from the 2017--18, 2018--19 and 2019--20
Premier League seasons, using the dataset's official split: 182 training, 35 validation and 38 test matches,
with no match in two splits. Each event is described by one frame per second from 30\,s before to 30\,s
after its timestamp. A frame-timing check found 6 of the 510 halves whose video timestamps drift more than
0.5\,s from the expected frame times; their 203 events are left out. After removing four non-event labels
(718 events), 22 event types remain with 15,860 events: 11,243 training, 2,307 validation and 2,310 test
(Table~\ref{tab:labels}).

\begin{table}[h]
\centering\footnotesize
\begin{tabular}{lrrrr@{\hspace{2.5em}}lrrrr}
\toprule
Kept label & Train & Val & Test & Total & Kept label & Train & Val & Test & Total\\
\midrule
corner & 1919 & 399 & 369 & 2687 & end\_of\_half\_game & 350 & 69 & 73 & 492\\
clearance & 1676 & 331 & 346 & 2353 & injury & 277 & 67 & 75 & 419\\
substitution & 997 & 192 & 211 & 1400 & free\_kick & 157 & 48 & 34 & 239\\
foul\_with\_no\_card & 1006 & 197 & 192 & 1395 & ball\_out\_of\_play & 168 & 37 & 30 & 235\\
shot\_off\_target & 906 & 162 & 164 & 1232 & var & 71 & 12 & 25 & 108\\
lead\_to\_corner & 775 & 177 & 163 & 1115 & foul\_lead\_to\_penalty & 68 & 14 & 15 & 97\\
goal & 624 & 125 & 126 & 875 & penalty & 39 & 8 & 9 & 56\\
yellow\_card & 555 & 126 & 118 & 799 & throw\_in & 28 & 14 & 7 & 49\\
saved\_by\_goal-keeper & 489 & 101 & 112 & 702 & red\_card & 20 & 5 & 3 & 28\\
off\_side & 386 & 87 & 71 & 544 & second\_yellow\_card & 8 & 3 & 0 & 11\\
ball\_possession & 365 & 63 & 92 & 520 & & & & & \\
start\_of\_half\_game & 359 & 70 & 75 & 504 & & & & & \\
\midrule
\textbf{Excluded label} & & & & & \textbf{Reason} & & & & \\
show\_added\_time & 325 & 68 & 73 & 466 & broadcast graphic & & & & \\
statistics\_and\_summary & 149 & 29 & 41 & 219 & broadcast graphic & & & & \\
unknown & 20 & 1 & 3 & 24 & scraper artefact & & & & \\
wcl-icon-settings-info-rounded & 9 & 0 & 0 & 9 & scraper artefact & & & & \\
\bottomrule
\end{tabular}
\caption{Event counts per label in the official split of the 255 matches, after leaving out the 6 halves that
failed the frame-timing check.}
\label{tab:labels}
\end{table}

\subsection*{2. Baseline results on real data}
Each frame is encoded with a frozen image or video model, and classifiers are trained on top. We compare two
frozen encoders: DINOv2-base~[3], a general-purpose image model, and MatchVision~[5], a video model
pretrained on soccer broadcasts that also attends across 30 frames. Every choice (pooling, classifier,
stopping epoch, model) is made on the 35 validation matches; the 38 test matches are scored once per seed.

\begin{table}[h]
\centering\footnotesize
\begin{tabular}{llcccc}
\toprule
Encoder & Model & Macro-F1 & Balanced acc. & Accuracy & 95\% range (acc.)\\
\midrule
-- & Majority class & -- & -- & $0.160$ & --\\
DINOv2 & Pooled features & $0.278 \pm 0.002$ & $0.273 \pm 0.007$ & $0.410 \pm 0.008$ & 0.389--0.432\\
DINOv2 & Sequence model & $0.390 \pm 0.009$ & $0.393 \pm 0.013$ & $0.509 \pm 0.005$ & 0.486--0.533\\
MatchVision & Pooled features & $0.385 \pm 0.019$ & $0.378 \pm 0.016$ & $0.510 \pm 0.012$ & 0.491--0.531\\
MatchVision & Sequence model & $\mathbf{0.405 \pm 0.001}$ & $\mathbf{0.406 \pm 0.008}$ & $\mathbf{0.523 \pm 0.007}$ & 0.502--0.546\\
\bottomrule
\end{tabular}
\caption{Test results on the 38 held-out matches, mean $\pm$ sd over 3 seeds. The majority class always predicts
the most common training label. The 95\% range comes from resampling whole test matches 1,000 times. These models were chosen by validation accuracy; the rerun with
selection by validation macro-F1 is in progress (Section~5).}
\label{tab:results}
\end{table}

\begin{itemize}[noitemsep]
  \item \textbf{Pooled frozen features}: the 60 frame vectors of an event pooled into one (mean, max, first or
        last frame, or the contrastive projection below), fed to logistic regression, a linear SVM or a small
        MLP; the best of these 15 pairings on validation is tested.
  \item \textbf{Sequence model}: four temporal models over the 60 frames (mean-pool MLP, BiLSTM, temporal CNN,
        Transformer) compared on validation. With DINOv2 it leads the pooled features by $+0.099$ accuracy
        (all 3 seeds); with MatchVision, which already models time, by only $+0.012$ (2 of 3 seeds).
  \item \textbf{Contrastive head}: trained with a triplet loss~[6] on training matches only, it raises event
        agreement of the clusters on held-out matches (ARI) from 0.039 to 0.109 with DINOv2 and from 0.094 to
        0.165 with MatchVision, while the agreement with match identity falls to 0.007 and 0.002.
  \item \textbf{Earlier data}, for comparison: on the 65 matches of 2021--22, 5-fold cross-validation by match
        gives macro-F1 $0.318 \pm 0.014$ and balanced accuracy $0.321 \pm 0.014$ for the sequence model with
        DINOv2, against $0.239 \pm 0.014$ and $0.243 \pm 0.011$ for pooled features (majority baseline 0.164).
\end{itemize}

\subsection*{3. Results per event type and the rare classes}
Following your advice, we kept all 22 types separate and looked at counts and per-class performance first
(Table~\ref{tab:perclass}).

\begin{table}[h]
\centering\footnotesize
\begin{tabular}{lrrcc@{\hspace{2.5em}}lrrcc}
\toprule
Event type & Train & Test & DINOv2 & MatchV. & Event type & Train & Test & DINOv2 & MatchV.\\
\midrule
corner & 1919 & 369 & 0.58 & 0.57 & ball\_possession & 365 & 92 & 0.19 & 0.24\\
clearance & 1676 & 346 & 0.49 & 0.49 & start\_of\_half\_game & 359 & 75 & 0.92 & 0.94\\
foul\_with\_no\_card & 1006 & 192 & 0.42 & 0.44 & end\_of\_half\_game & 350 & 73 & 0.71 & 0.68\\
substitution & 997 & 211 & 0.74 & 0.77 & injury & 277 & 75 & 0.54 & 0.40\\
shot\_off\_target & 906 & 164 & 0.38 & 0.43 & ball\_out\_of\_play & 168 & 30 & 0.10 & 0.02\\
lead\_to\_corner & 775 & 163 & 0.32 & 0.39 & free\_kick & 157 & 34 & 0.22 & 0.13\\
goal & 624 & 126 & 0.67 & 0.67 & var & 71 & 25 & 0.45 & 0.42\\
yellow\_card & 555 & 118 & 0.50 & 0.62 & foul\_lead\_to\_penalty & 68 & 15 & 0.29 & 0.27\\
saved\_by\_goal-keeper & 489 & 112 & 0.22 & 0.33 & penalty & 39 & 9 & 0.06 & 0.26\\
off\_side & 386 & 71 & 0.38 & 0.46 & throw\_in & 28 & 7 & 0.00 & 0.00\\
 & & & & & red\_card & 20 & 3 & 0.00 & 0.00\\
 & & & & & second\_yellow\_card & 8 & 0 & -- & --\\
\bottomrule
\end{tabular}
\caption{Test F1 per event type of the sequence model, mean over 3 seeds, for each encoder, with training and
test counts. \texttt{second\_yellow\_card} has no test events.}
\label{tab:perclass}
\end{table}

\begin{itemize}[noitemsep]
  \item \textbf{Too rare to learn.} The three types with fewer than 30 training events are not learned:
        \texttt{red\_card} (20) and \texttt{throw\_in} (28) score F1 = 0 with both encoders, and
        \texttt{second\_yellow\_card} (8) has no test events at all.
  \item \textbf{Rarity is not the whole story.} \texttt{ball\_out\_of\_play} (168 training events) and
        \texttt{ball\_possession} (365) stay below 0.25. They look like ordinary play on screen, so more
        examples alone may not fix them.
  \item \textbf{Proposal, not applied.} Merge \texttt{second\_yellow\_card} into \texttt{red\_card} as one
        ``sending off'' type: a second yellow card is a sending off, so the broader category makes sense, and
        together they still have only 28 training events. We see no meaningful broader category for
        \texttt{throw\_in}, so we would keep it and report it as not yet learnable.
\end{itemize}

\subsection*{4. Season hold-out: 2022--23 reserved}
Table~\ref{tab:seasons} lists the Premier League matches available. 2022--23 is the most recent season with
enough matches for a test set and has never been used, so we have reserved all 115 of its matches: they are
listed in a manifest marked as test, and nothing reads their frames, features or labels until our final
models are fixed. We will train and validate on 2017--18 to 2021--22 (411 matches). 2023--24 is more recent
but has only 26 matches. 2021--22 cannot be the hold-out, since our design decisions were made on it.

\begin{table}[h]
\centering\small
\begin{tabular}{lccc}
\toprule
Season & Matches & Official train / val / test & Used in development\\
\midrule
2017--18 to 2019--20 & 255 & 182 / 35 / 38 & yes (this report)\\
2020--21 & 91 & 68 / 13 / 10 & no\\
2021--22 & 65 & 50 / 7 / 8 & yes (weeks 2--5)\\
\textbf{2022--23} & \textbf{115} & \textbf{89 / 13 / 13} & \textbf{no: reserved for the final test}\\
2023--24 & 26 & 21 / 3 / 2 & no\\
\bottomrule
\end{tabular}
\caption{Premier League matches in SoccerReplay-1988, with the dataset's official split.}
\label{tab:seasons}
\end{table}

\subsection*{5. Model selection by macro-F1}
Until this week our classifiers were chosen by validation accuracy, although macro-F1 is our main metric, and
the two can disagree. On the MatchVision pooled features, for example, the best candidate on accuracy
(max-pool with an MLP) has validation macro-F1 0.347, while the best on macro-F1 (mean-pool with a linear
SVM) has 0.363. Every classifier now ranks its candidates by validation macro-F1, with accuracy breaking ties,
and the sequence models also stop training on it. Test matches are still used only once. The rerun on the
255 matches is in progress; Table~\ref{tab:results} will be updated with its results.

\section*{Response to Last Week's Comments}
\begin{enumerate}[noitemsep]
  \item \textbf{Keep the rare classes separate; look at counts and per-class performance first.} Done: no
        classes were merged; Tables~\ref{tab:labels} and~\ref{tab:perclass} give the counts and per-class F1.
        Only one merge passes both of your conditions (too rare to learn, and a sensible broader category):
        \texttt{second\_yellow\_card} into \texttt{red\_card}. We ask before applying it (see Questions).
  \item \textbf{Hold out the most recent unused season, untouched until the final test.} Done: 2022--23,
        115 matches, reserved (Section~4).
  \item \textbf{Baseline results on real data with macro-F1 and balanced accuracy.} Done:
        Table~\ref{tab:results}, on 255 matches, with the majority baseline and intervals that resample
        whole test matches. Model selection now also uses macro-F1.
\end{enumerate}

\section*{Challenges and Issues}
\textbf{MatchVision and the hold-out season.} MatchVision was pretrained on SoccerReplay-1988's official
training matches, so 89 of the 115 reserved matches may have been in its pretraining data, which would make
its hold-out score optimistic. DINOv2 was not trained on this dataset. We will report MatchVision both on all
115 matches and on the 26 official validation and test matches of 2022--23, which it did not train on.

\textbf{Annotations for the hold-out.} The annotation files of the dataset's training archive are not yet
installed for 2022--23, so 89 of the 115 reserved matches cannot be extracted yet. This does not delay the
final test, which comes last.

\textbf{Visually ambiguous types.} Several common types (\texttt{ball\_out\_of\_play},
\texttt{ball\_possession}, \texttt{free\_kick}) score low despite hundreds of examples, which limits macro-F1
whatever we do about the rare classes.

\section*{Next Steps}
\begin{enumerate}[noitemsep]
  \item Finish the macro-F1 selection rerun and update Table~\ref{tab:results}.
  \item Apply your decision on the rare classes and rerun once.
  \item Extract 2020--21 and 2021--22 together with the 255 matches, so training and validation use 411
        matches, with five seeds.
  \item Keep 2022--23 untouched until the final models are fixed, then test once.
\end{enumerate}

\section*{Questions or Feedback Needed}
\begin{itemize}[noitemsep]
  \item May we merge \texttt{second\_yellow\_card} into \texttt{red\_card} as one ``sending off'' type, and keep
        \texttt{throw\_in} separate, reported as not yet learnable?
  \item For MatchVision on the hold-out season, is reporting both all 115 matches and the 26 matches outside
        its pretraining data an acceptable way to handle the overlap?
\end{itemize}

\section*{Team Contributions}
\begin{itemize}[noitemsep]
  \item \textbf{Aniyah Mack:} [to be filled in by Aniyah]
  \item \textbf{Adonay Zeru:} [to be filled in by Adonay]
  \item \textbf{Hasset Admasu:} Extended the data to 255 matches on the server, ran the baseline experiments
        with both encoders, computed the per-class results, chose and reserved the hold-out season, and
        switched model selection to macro-F1, working with Claude Code (see below).
  \item \textbf{Nathnael Abraham:} [to be filled in by Nathnael]
\end{itemize}

\section*{References}
\begin{enumerate}[label={[\arabic*]}, noitemsep]
  \item Claude Code (AI coding assistant). Wrote most of the pipeline and analysis code and the data checks
        under our direction; we chose what to build and what to keep, ran it on the server, and reviewed
        the results.
  \item James MacQueen. Some methods for classification and analysis of multivariate observations. In
        \emph{Proceedings of the Fifth Berkeley Symposium on Mathematical Statistics and Probability, Volume 1:
        Statistics}, pages 281--297. University of California Press, 1967.
  \item Maxime Oquab, Timoth\'ee Darcet, Th\'eo Moutakanni, Huy V. Vo, Marc Szafraniec, Vasil Khalidov, Pierre
        Fernandez, Daniel Haziza, Francisco Massa, Alaaeldin El-Nouby, Mahmoud Assran, Nicolas Ballas, Wojciech
        Galuba, Russell Howes, Po-Yao Huang, Shang-Wen Li, Ishan Misra, Michael Rabbat, Vasu Sharma, Gabriel
        Synnaeve, Hu Xu, Herv\'e J\'egou, Julien Mairal, Patrick Labatut, Armand Joulin, and Piotr Bojanowski.
        DINOv2: Learning robust visual features without supervision. \emph{Transactions on Machine Learning
        Research}, 2024.
  \item Alec Radford, Jong Wook Kim, Chris Hallacy, Aditya Ramesh, Gabriel Goh, Sandhini Agarwal, Girish Sastry,
        Amanda Askell, Pamela Mishkin, Jack Clark, Gretchen Krueger, and Ilya Sutskever. Learning transferable
        visual models from natural language supervision. In \emph{Proceedings of the 38th International
        Conference on Machine Learning (ICML)}, volume 139, pages 8748--8763. PMLR, 2021.
  \item Jiayuan Rao, Haoning Wu, Hao Jiang, Ya Zhang, Yanfeng Wang, and Weidi Xie. Towards universal soccer
        video understanding. In \emph{Proceedings of the IEEE/CVF Conference on Computer Vision and Pattern
        Recognition (CVPR)}, pages 8384--8394, 2025.
  \item Florian Schroff, Dmitry Kalenichenko, and James Philbin. FaceNet: A unified embedding for face
        recognition and clustering. In \emph{Proceedings of the IEEE Conference on Computer Vision and Pattern
        Recognition (CVPR)}, pages 815--823, 2015.
  \item Aniyah Mack, Adonay Zeru, Hasset Admasu, and Nathnael Abraham. Soccer Event Baseline: supervised
        mean-pooled baseline for soccer event recognition, with weekly check-ins. GitHub repository,
        \url{https://github.com/Hasset-hub/soccer-event-baseline}, 2026.
  \item Laurens van der Maaten and Geoffrey Hinton. Visualizing data using t-SNE. \emph{Journal of Machine
        Learning Research}, 9(86):2579--2605, 2008.
\end{enumerate}

\end{document}
